\geometry{paperwidth=22cm, paperheight=8cm, margin=0.5cm}
\begin{document}

% Seitengeometrie für Tabu-Karten (Querformat 21:9)

\pagestyle{empty}

\huge
\newcommand{\tabucard}[2]{
    \begin{tcolorbox}[
        colback=white,
        colframe=black,
        boxrule=2pt,
        arc=0pt,
        height=\textheight,
        width=\textwidth,
        valign=center,
        halign=center,
        nobeforeafter
    ]
        \vspace*{\fill}
        \begin{center}
            \textbf{\Huge #1}
            
            \vspace{1.5cm}
            
            \begin{tcolorbox}[
                colback=lightgray,
                colframe=black,
                boxrule=1.5pt,
                arc=4pt,
                width=0.85\textwidth,
                nobeforeafter
            ]
                \textbf{\Large Tabu-Wörter:} \\[0.5cm]
                \Large #2
            \end{tcolorbox}
        \end{center}
        \vspace*{\fill}
    \end{tcolorbox}
    \newpage
}

\newgeometry{paperwidth=42cm,paperheight=18cm,margin=0.5cm}

\section*{\Huge Tabu}
\textbf{Ablauf}
\begin{itemize}
\item \textbf{2 Gruppen}
\item \textbf{Ziel}: Möglichst viele Wörter erraten
\item \emoji{hourglass} \textbf{Zeit}: \textbf{1 Minute} pro Gruppe 
\item Möglichkeit, \textbf{1 Wort pro Runde} zu \textbf{überspringen}
\item Je nach Zeit: \textbf{Finale} \emoji{trophy}
\end{itemize}

\textbf{Regeln}:
\begin{itemize}
\item \textbf{Verboten}: Wort selber und Tabu-Wörter
\item \textbf{Keine Wörter übersetzen} (weder Wort selber noch Tabu-Wörter)
\item \textbf{Keine Teil-Wörter} erlaubt (z.B. Fahrradlenker $\rightarrow$ \sout{Fahrrad}, \sout{Lenker})
\end{itemize}

\newpage
\restoregeometry

\tabucard{Beispiel: Turtle}{
Schildkröte, zeichnen, Python, Code, Spirale
}

